\section{Language}
\label{sec:formal}

\begin{figure}[t!]
  \vspace{-.3cm}
{\footnotesize
  \begin{alignat*}{2}
    \text{\textbf{Variables }}& \quad &\quad& x, y, z, \vnu, ...
    \\\text{\textbf{Base Types}}& \quad & b  ::= \quad &  \Code{unit} ~|~ \Code{bool} ~|~ \Code{nat} ~|~ \Code{int} ~|~ \Code{id} ...
    \\\text{\textbf{Basic Types}}& \quad & s  ::= \quad &  b ~|~ s\sarr s
    \\\text{\textbf{Pure Operations}}& \quad & \primop ::= \quad & {+} ~|~ {-} ~|~ {==} ~|~ {<} ~|~ {\leq} ~|~ ...
    \\ \text{\textbf{Constants }}& \quad &c ::= \quad & () ~|~ \mathbb{B} ~|~ \mathbb{Z} ~|~ ...
    \\ \text{\textbf{Qualifiers}}& \quad & \phi ::= \quad & c ~|~ x ~|~ \primop\;\overline{v} ~|~ \bot ~|~ \top  ~|~ \neg \phi ~|~ \phi \land \phi ~|~ \phi \lor \phi ~|~ \phi \impl \phi ~|~ \forall x{:}b.\, \phi
    \\\text{\textbf{Effectful Operations}}& \quad & \eff{op} ::= \quad & \eff{read} ~|~ \eff{write} ~|~ ...
    % \\\text{\textbf{Event Kinds}}& \quad & k ::= \quad & \eff{gen}
    % ~|~ \eff{obs}
    \\ \text{\textbf{Events}} & \quad & m ::= \quad & \eff{op}{\; \overline{v}} ~|~ \reqOp{op}{\ensuremath{\langle \iota,\overline{v}\rangle}} ~|~ \respOpAlt{op}{\ensuremath{\iota}}
    \\ \text{\textbf{Values }}& \quad & v ::= \quad & c ~|~ x ~|~ \zlam{x}{s}{e} ~|~ \zfix{f}{s}{x}{s}{e}
    \\ \text{\textbf{Expressions}}& \quad & e ::=\quad & v ~|~ e
    \oplus e ~|~ v\; v ~|~ \primop\; \overline{v} ~|~ m
    ~|~\zlet{x{:}b}{e}{e} ~|~ \sassume{\phi} ~|~ \sassert{\phi}
    \\[5pt]
    \text{\textbf{Traces}} & \quad & \alpha ::= \quad & \emptr ~|~ m
    \;{::}\; \alpha ~|~ \alpha \listconcat \alpha
    \\[5pt]
    \text{\textbf{Notations}} & \quad & e_1;e_2 \doteq \quad & \zlet{\_}{e_1}{e_2}
    %  \mykw{let}\;\eff{op}\;\overline{v}\;\mykw{in}\;e \doteq \zgen{op}{\overline{v}}{e} \qquad\qquad
  \end{alignat*}
}
\vspace*{-.25in}
\caption{Syntax of \DSL{} expressions}
\label{fig:term-syntax}
\vspace*{-.1in}
\end{figure}

We formalize our approach using a core language, \DSL{}, for
expressing \textit{trace generators}-- nondeterministic programs that
probe the behaviors of a black-box system under test via a set of
(effectful) operators that the system provides. This calculus is
equipped with a type system that characterizes the set of traces that
a program \emph{must} be able to produce, \emph{if} the system under
test produces appropriate responses. The syntax of \DSL{} is shown in
\autoref{fig:term-syntax}; the calculus includes both pure and
effectful operations ($\primop$ and $\eff{op}$), nondeterministic
choices ($\oplus$), and $\mykw{assume}$ and $\mykw{assert}$ statements
in the style of solver-aided languages~\cite{BC+10}. Effectful
operations can be tagged as a request (\req{}), indicating that it
should be handled asynchronously; performing this operation generates
an id (given as $\iota$ in the syntax) that can then be subsequently
used when performing a corresponding response operation (\resp{}).
%\autoref{fig:term-syntax} defines some shorthand notations used by our
%examples, e.g., we use $e_1;e_2$ to simplify $\mykw{let}$-expressions
%when the remaining program does not use the bound variable, i.e.,
%$\zlet{\_}{e_1}{e_2}$.
\DSL{} is equipped with a similar small-step operational semantics as
other calculi for trace-based type systems~\cite{ZYDJ23}.\footnote{The
  reduction rules are included in the appendix.} Like those systems,
the reduction relation $\steptr{\alpha}{e}{~\alpha'}{e'}$, includes a
\textit{history trace} $\alpha$, that constrains the result of an
effectful operation-- for example, the result of a $\eff{get}$ effect
depends on the $\eff{put}$ operations that preceded it.

% \begin{example}[Auxiliary Operator Semantics]\label{ex:aux-semantics}
% The auxiliary semantics of the $\eff{put}$ and $\eff{get}$ operators in Example~\ref{ex:sc} can be defined as follows:

% {\footnotesize
%   \mprooftr{42mm}{}{OpPut}{
%   $\alpha \vDash \zevent{putReq}{k,v} \Downarrow \{ \zevent{putRsp}{k,v} \}$
%   }
%   \quad
%   \mprooftr{62mm}{}{OpGet}{
%   $\alpha_1\listconcat\zevent{putRsp}{k,v}\listconcat\alpha_2\vDash \zevent{getReq}{k} \Downarrow \{\zevent{getRsp}{k,v}\}$
%   }
%   \\[-0.1em]
%   }\noindent

% \noindent
% The database handles the $\eff{putReq}$ event by simply adding a $\eff{putRsp}$ event with the same payload to the capability. To guarantee the EC, handling a $\eff{getReq}$ adds a $\eff{getRsp}$ event where the read value is the one previously written value $v$ at the time the $\eff{getReq}$ occurs.
% \end{example}

% \begin{example}[Operational Semantics]\label{ex:semantics-stlc}
% We show a program that generates STLC term $\lambda.[0]$.
%   {\footnotesize
%     \begin{flalign*}
%       &[] \vDash
%       (\emptyset, \zassume{x}{\top}{\eff{abs}\;x;\eff{var}\;0;\mykw{obs}\;\eff{endAbs}})
%       \xhookrightarrow{[]} (\emptyset, \eff{abs}\;\Int;\eff{var}\;0;\mykw{obs}\;\eff{endAbs}) \tag{\textsc{StAssume}} &&
%       \\[-0.4em]
%       &\;\;\xhookrightarrow{[\zevent{abs}{\Int}]} (\{\eff{endAbs}\}, \eff{var}\;0;\mykw{obs}\;\eff{endAbs})  \xhookrightarrow{[\zevent{var}{0}]} (\{\eff{endAbs}\}, \mykw{obs}\;\eff{endAbs})
%       \xhookrightarrow{[\eff{endAbs}]} (\emptyset, ()) \tag{\textsc{StGen, StObs}} &&
%     \end{flalign*}
%   }\noindent
%     The program first instantiates $x$ with an arbitrary STLC type (e.g., $\Int$), applies $\eff{abs}$ to produce capability $\eff{endAbs}$, executes $\eff{var}$ which generates no new capability, and finally consumes $\eff{endAbs}$.
% \end{example}

% \begin{example}[Strong consistency violation]\label{ex:semantics-async}
%   The following program shows a strong consistency violation in Example~\ref{ex:sc}, under the auxiliary semantics of $\eff{put}$ and $\eff{get}$ shown in Example~\ref{ex:aux-semantics}.
%     {\footnotesize
%       \begin{flalign*}
%         &[\zevent{putReq}{``\Code{k}", 1};\zevent{putRsp}{``\Code{k}", 1}] \vDash
%         (\emptyset, \eff{putReq}\;``\Code{k}"\;2;\eff{getReq}\;``\Code{k}";\sobs{putRsp};\zobs{k, v}{getRsp}{\sassert{v \not= 2}})
%         \\[-0.2em]
%         &\;\;\xhookrightarrow{[\zevent{putReq}{``\Code{k}", 2}]} (\{ \zevent{putRsp}{``\Code{k}", 2} \}, \eff{getReq}\;``\Code{k}";\sobs{putRsp};\zobs{k, v}{getRsp}{\sassert{v \not= 2}})  \tag{\textsc{StGen}} &&
%         \\[-0.4em]
%         &\;\;\xhookrightarrow{[\zevent{getReq}{``\Code{k}"}]} (\{ \zevent{putRsp}{``\Code{k}", 2}, \zevent{getRsp}{``\Code{k}", 1} \}, \sobs{putRsp};\zobs{k, v}{getRsp}{\sassert{v \not= 2}})  \tag{\textsc{StGen}} &&
%         \\[-0.4em]
%         &\;\;\xhookrightarrow{[\zevent{putRsp}{``\Code{k}", 2}]} (\{\zevent{getRsp}{``\Code{k}", 1}\}, \zobs{k, v}{getRsp}{\sassert{v \not= 2}})  \tag{\textsc{StObs}} &&
%         \\[-0.4em]
%         &\;\;\xhookrightarrow{[\zevent{getRsp}{``\Code{k}", 1}]} (\emptyset, \sassert{1 \not= 2})  \tag{\textsc{StObs}} &&
%       \end{flalign*}
%     }\noindent
%   As shown above, the program update the value of key $\Code{k}$ to $2$, but the read value is still $1$ due to the auxiliary semantics of $\eff{getReq}$ with EC is allowed to return any previous written value.
% \end{example}

% We show a program that generates the identity function $\lambda.[0]$ in the style of \autoref{fig:spec} as follows:
% {\footnotesize
%   \begin{flalign*}
%     &[] \vDash
%     (\emptyset, \zassume{x}{\top}{\eff{abs}\;x;\eff{var}\;0;\eff{putTy}\;x;\eff{ayncGetTy};\zlet{y}{\mykw{obs}\;\eff{getTy}}{\mykw{obs}\;\eff{endAbs};\eff{putTy}\;\Code{Arr}(x, y)}})\\[-0.4em]
%     &\;\;\xhookrightarrow{[]} (\emptyset, \eff{abs}\;\Int;\eff{var}\;0;\eff{putTy}\;\Int;\eff{ayncGetTy};\zlet{y}{\mykw{obs}\;\eff{getTy}}{\mykw{obs}\;\eff{endAbs};\eff{putTy}\;\Code{Arr}(\Int, y)}) \tag{\textsc{StAssume}} &&
%     \\[-0.4em]
%     &\;\;\xhookrightarrow{[\zevent{abs}{\Int}]} (\{\eff{endAbs}\}, \eff{var}\;0;\eff{putTy}\;\Int;\eff{ayncGetTy};\zlet{y}{\mykw{obs}\;\eff{getTy}}{\mykw{obs}\;\eff{endAbs};\eff{putTy}\;\Code{Arr}(\Int, y)})  \tag{\textsc{StGen}} &&
%     \\[-0.6em]
%     &\;\;\xhookrightarrow{[\zevent{var}{0};\zevent{putTy}{\Int};\eff{ayncGetTy}]}^* (\{\eff{endAbs};\zevent{getTy}{\Int}\}, \zlet{y}{\mykw{obs}\;\eff{getTy}}{\mykw{obs}\;\eff{endAbs};\eff{putTy}\;\Code{Arr}(\Int, y)}) \tag{\textsc{StGen}} &&
%     \\[-0.6em]
%     &\;\;\xhookrightarrow{[\zevent{getTy}{\Int}]} (\{\eff{endAbs}\}, \mykw{obs}\;\eff{endAbs};\eff{putTy}\;\Code{Arr}(\Int, \Int)) \tag{\textsc{StObs}} &&
%     \\[-0.4em]
%     &\;\;\xhookrightarrow{[\eff{endAbs}]} (\emptyset, \eff{putTy}\;\Code{Arr}(\Int, \Int)) \xhookrightarrow{[\zevent{putTy}{\Code{Arr}(\Int, \Int)}]} (\emptyset, ()) \tag{\textsc{StObs, StGen}} &&
%   \end{flalign*}
% }\noindent
%   The first step substitutes $x$ with an arbitrary STLC type (e.g., $\Int$). In the next steps, the program performs the generative operator $\eff{abs}$, which generates a new capability $\eff{endAbs}$ to be handled in the future. Similarly, the program performs $\eff{var}$, $\eff{putTy}$, and $\eff{ayncGetTy}$, where only the last operator generates a new capability, $\zevent{getTy}{\Int}$, whose payload is the last $\eff{put}$ type.
%   The fourth step consumes the $\zevent{getTy}{\Int}$ event and substitutes $y$ with $\Int$ in the remaining program. The rest of the program is then executed, which ends the lambda abstraction and records the type of the generated STLC term as $\Code{Arr}(\Int, \Int)$.

\subsection{Type System}

\begin{figure}[t!]
{\small
\begin{alignat*}{2}
    % \\\text{\textbf{Symbolic Event}}& \quad &se  ::= \quad & \msgB{op}{\overline{x}}{\phi}
    \text{\textbf{Symbolic Regex}}& \quad &H,F,A  ::= \quad & \emptyset ~|~ \epsilon ~|~ \msgB{op}{\overline{x}}{\phi} ~|~ A \lorA A ~|~ A \seqA A ~|~ A^* ~|~ A \setminus A ~|~ A \landA A ~|~ \anyA
    % \\& \quad &  \quad & \text{\textcolor{gray}{in Symbolic LTL$_f$ }}  se ~|~ A \lorA A ~|~ \nextA A ~|~ A \untilA A ~|~ \negA A ~|~ A \landA A
    \\\text{\textbf{Pure Refinement Types}}& \quad & t  ::= \quad & \urt{b}{\phi} ~|~  x{:}t\sarr t ~|~ x{:}t\sarr \tau
    \\ \text{\textbf{Underapproximate HATs}}& \quad & \tau  ::= \quad & \uhat{H}{x{:}t}{F} ~|~ \tau \interty \tau ~
    |~ x{:}b \garr \tau ~|~ \effGray{op}{:}s \sarr \tau
    \\\text{\textbf{Type Contexts}}& \quad &\Gamma ::= \quad  &\emptyset ~|~ x{:}t, \Gamma
  \end{alignat*}
  \vspace*{-.15in}
  {\footnotesize\begin{flalign*}
    {\small\text{\textbf{Type Aliases}}} & \;\; \msg{op}{\phi} \doteq \msgB{op}{\overline{x}}{\phi} \;\; \msgA{op}{\overline{v}} \doteq \msgB{op}{\overline{x}}{\overline{x = v}} \;\; \msgC{op} \doteq \msgB{op}{\overline{x}}{\top} \;\; A_1 A_2 \doteq A_1 \seqA A_2 \;\; b \doteq \urt{b}{\top}
    % \;\; \tau \doteq x{:}b\garr\tau
    % \\
    % & \widehat{\ort{b}{\phi}} \doteq \urt{b}{\phi} \;\; \widehat{\urt{b}{\phi}} \doteq \ort{b}{\phi} \;\; \widehat{x{:}t\sarr t'} \doteq x{:}t\sarr t' \;\; \widehat{x{:}t\sarr \tau} \doteq x{:}t\sarr \tau
  \end{flalign*}}
}
\vspace*{-.15in}
\caption{\DSL{} types}
\label{fig:type-syntax}
\vspace*{-.15in}
\end{figure}

The syntax of types in \DSL{} is shown in \autoref{fig:type-syntax}.
Types include \emph{pure} refinement types, which describe pure
computations, and underapproximate Hoare Automata Types (\tyName{}s),
which describe effectful computations.
% which describe incorrectness triple $\effLR{H}\;e\;\effLR{F}$ as a
% type $\uhat{H}{\Unit}{F}$ of term $e$.
Pure refinement types are similar to those found in prior
underapproximate refinement type systems~\cite{CoverageType}, and
allow base types (e.g., {\Code{int}}) to be further constrained by a
logical formula or qualifier. 

In contrast to traditional type systems where pure type qualifier $\ort{b}{\phi}$ 
constrains every value a pure expresion \emph{may} evaluate to; 
the type qualifier $\urt{b}{\phi}$ constrains a subset of the values it
\emph{must} evaluate to.  
For example, the constant value $1$ can be
assigned the types $\ort{\Int}{\vnu = 1}$ and $\ort{\Int}{0 < \vnu}$
-- it satisfies constraint in both \emph{overapproximate} qualifiers
-- but, it cannot be assigned the type, $\urt{\Int}{\vnu = 1 \lor \vnu = 2}$, since the value 1 cannot evaluate to both 1 or 2 as required
by this \emph{underapproximate} qualifier. The term $1 \oplus 2 \oplus
3$, in contrast, does have this underapproximate type.

\paragraph{Underapproximate Hoare Automata Types}

Similar to other recent trace-based type systems~\cite{ZYDJ24}, \DSL{}
uses \emph{Symbolic Finite Automata} (SFAs)~\cite{Vea13,
  SFA-minimization, SFA-Transducers} to qualify the traces generated
by effectful computations; in contrast to those prior works, the
\tyName{}s in $\DSL{}$ use SFAs to \emph{underapproximate} the set of effects
that a program must produce. A \tyName{} $\uhat{H}{x{:}t}{F}$ consists
of a pure refinement type $t$ and two SFAs: a \emph{history} SFA $H$
that constrain the history trace under which a term is executed, and a
\emph{future} SFA $A$ that describes a subset of the events that it
must generate as a result. \tyName{} represents SFAs using symbolic
regular expressions (SREs), but in principle could support any
encoding of SFAs, e.g., Symbolic LTL$_f$~\cite{LTLf}. SREs support
\emph{symbolic events}, $\msgB{op}{\overline{x}}{\phi}$, atomic
predicates that describes an effectful operation {$\eff{op}$} whose
arguments $\overline{x}$ satisfy the qualifier $\phi$. SREs provide a
convenient language for writing SFAs, supporting versions of the
standard regex operators shown in \autoref{fig:type-syntax}, including
union $A \lorA A$, sequence $A \seqA A$, and kleene star $A^*$
operators.

\subsubsection{Example}\label{ex:stack-example}
  Consider an effectful stack library implemented as a fixed-size
  array whose elements grow monotonically.  An incorrect implementation
  that fails to grow the array when the stack is full may cause
  $\eff{push}$ operations to be dropped.  To detect this violation, we
  need to explore traces of the following form, where the last pop
  operation raises an ``empty stack'' exception:
  \begin{align*}
    \zevent{push}{1};\zevent{push}{2}; \zevent{push}{3}; \zevent{pop}{3}; \zevent{pop}{2};\textcolor{red}{\zevent{pop}{1}}
  \end{align*}\noindent
  Because we would like the \tyName{} specifications for $\eff{push}$
  and $\eff{pop}$ to guide the synthesis of a test generator, they
  should be interpreted as underapproximations, consisting of nested
  pairs of $\eff{push}$ and $\eff{pop}$ events. This can be expressed
  by the following \tyName{} specifications, using ghost events
  $\effGray{pushI}$ and $\effGray{popI}$ as shown below:
  {\small\begin{align*} &\eff{push} :~ y{:}\Int \garr
    x{:}\urt{\Int}{y < \vnu} \sarr
    \uhat{\Code{Last}(\msgAGray{pushI}{n\;y})}{\Unit}{\msgA{push}{x}\msgAGray{pushI}{n{+}1\;x}}
    \\&\eff{pop} :~ \uhat{\Code{Last}(\msgAGray{popI}{m\;\_}) \land
      \Code{Last}(\msgAGray{pushI}{n\;\_}) \land
      (\allA\msgAGray{pushI}{(n{-}m)\;x}\allA)}{x{:}\Int}{\msgA{pop}{x}\msgAGray{popI}{m{+}1\;
        x}}
  \end{align*}}\noindent
 Intuitively, the $\effGray{pushI}$ and $\effGray{popI}$ ghost events are
 simply $\eff{push}$ and $\eff{pop}$ events equipped with indices to
 record the number of push and pop operations\footnote{For simplicity,
   we assume the index starts from $1$, so the first push will be
   recorded as $\msgAGray{pushI}{1\;x}$. We also omit unused variables
   using the wildcard $\_$.}.  The signature for $\eff{push}$ requires
 that the new pushed value ($x$) is greater than the last pushed value
 ($y$), and adds a new ghost event $\effGray{pushI}$ with an
 incremented index (i.e., $\msgAGray{pushI}{n{+}1\;x}$).  The
 signature for $\eff{pop}$ also records the number of pop operations
 in the same way (i.e., $\msgAGray{popI}{m{+}1\;x}$); additionally, it
 requires the last popped value ($x$) to be equal to the
 $(n{-}m)^{th}$ pushed value, reflecting the ``last-in-first-out’'
 property of the stack.
%%\end{example}
%% \tyName{}s can only appear in the result type of a function, where
%% they describe the effects performed when the function is evaluated.
%% Following prior work~\cite{ZYDJ24}, our type system includes
%% \emph{ghost variables} ($x{:}b \garr \tau$), allowing it to explicitly
%% capture data dependencies between symbolic events in \tyName{}s. It
%% also includes analogous \emph{ghost events} ($\effGray{op}{:}s \sarr
%% \tau$, which are written in $\effGray{gray}$ to distinguish them from
%% regular effect operators), enabling \tyName{}s to describe non-regular
%% patterns of effects, e.g., the nesting depth of abstractions in an
%% STLC terms as described in the previous section.  Like ghost
%% variables, ghost operators are used exclusively for specification and
%% verification purposes--= they do not appear in expressions and do not
%% impact the evaluation of a program.

%% \SJ{This example is too detailed.  We should use something simpler -
%%   maybe a stack example that counts pushes and pops?}  \BD{Do we need
%%   to say anything more about ghost events than what we said in the
%%   overview -- our page budget is probably spent on additional examples
%%   of the novel features of the type system.}
%% \begin{example}[Ghost Events]\label{ex:stlc-types}
%%   The \tyName{} specification of the %
%%   $\eff{token}{:}\Code{str} \sarr \Unit$ %
%%   operator from ~\autoref{ex:de-bruijn-index} is:%
%%   \vspace{-.4cm} \par\nobreak {\footnotesize
%%     \begin{align*}
%%       & \effGray{depth}{:}(\Int\sarr\Unit) \garr d{:}\Int \garr
%%       \\&\;\;(t{:}\Code{stlcTy}\garr x{:}\ort{\Code{str}}{\vnu = \boxed{)}} \sarr \uhat{\Code{LAST}(\msgAGray{depth}{d}) \msgA{token}{\boxed{(\lambda t.}\;} \msgAGray{depth}{d + 1} \allA }{\Unit}{\msgA{token}{x} \msgAGray{depth}{d}}) \interty
%%       \\&\;\;(n{:}\Int\garr x{:}\ort{\Code{str}}{\vnu = \boxed{[n]} \land 0 \leq n < d} \sarr \uhat{\Code{LAST}(\msgAGray{depth}{d})}{\Unit}{\msgA{token}{x}})
%%     \end{align*}}\noindent
%%   This \tyName{} consists of two intersected \tyName{} for function
%%   abstraction and bound variable (parameter $x$ is constrained to be
%%   $\boxed{)}$ and $\boxed{[n]}$ respectively), which share the same
%%   ghost operator $\effGray{depth}$ and ghost variable $d$.
%%   \BD{This example needs to be gentler...}
%% \end{example}


% The standard temporal operators
% (e.g., test $\evparenth{\phi}$, next $\nextA A$, until $\untilA$) and
% various set operators (i.e., complement $\neg$, intersection $\land$,
% and union $\lorA$) are defined normally. These operators are
% sufficient to capture other modalities, e.g., eventually ($\finalA$),
% globally ($\globalA$), and importantly, the singleton (last) modality
% $$, which describes a singleton trace, i.e., one which prohibits
% any subsequent effects~\cite{LTLf}. SREs can capture several common
% patterns: the set of all possible traces \(\allTL\), the singleton set
% containing the empty trace \(\epsilon\), and the empty set of traces and
% \(\empTL\); these are analogous to the regular expressions \( .^* \),
% \(\epsilon\), and \(\emptyset\), resp.

% \begin{example}[Queue Library]\label{ex:stack}
% The precise semantics of $\eff{enqueue}$, $\eff{dequeue}$ operation of an effectful queue library requires \emph{first-in-first-out} (FIFO) property which needs counting. We introduce two ghost event $\effGray{numEnq}$ and $\effGray{numDeq}$ which are counters that records the number of previous added and removed elements. Thus, the FIFO property can be defined as ``the $i^{th}$ dequeued element is exactly the $i^{th}$ enqueued element'', and expressed as the following \tyName{}:
% {\footnotesize
%     \begin{align*}
%       \eff{enqueue} :~ & i{:}\Int\garr x{:}\Int\sarr\uhat{\Code{LAST}(\msgAGray{numEnq}{i - 1})}{\msgA{enqueue}{x}\msgAGray{numEnq}{i}}
%       \\\eff{dequeue} :~ & i{:}\Int\garr x{:}\Int\sarr\uhat{\Code{LAST}(\msgA{enqueue}{x}\msgAGray{numEnq}{i}) \land \Code{LAST}(\msgAGray{numDeq}{i - 1})}{\msgA{dequeue}{x}\msgAGray{numDeq}{i}}
%     \end{align*}}\noindent
%   The history regex of $\eff{enqueue}$ operator maintains the counter by get number of previous added elements by match the most last $\effGray{numEnq}$ event in the history (we assume this is number is $0$ when there is no previous $\effGray{numEnq}$) and increase this number by $1$ by appending a new ghost event in future regex. The \tyName of $\eff{dequeue}$ besides maintians the counter $\effGray{numDeq}$ in the same way, also requires the $\eff{dequeue}$ element is exactly the $i^th$ $\eff{enqueue}$ element, i.e., the one update the counter $\effGray{numDeq}$ as $i$.
% \end{example}

\autoref{fig:type-syntax} gives the definition of the type aliases
that \autoref{sec:overview} used to limit user annotation burden.
When the fields of an effectful operation are clear from context, we
omit its arguments $\overline{x}$, e.g., $\msg{\req(\textbf{write})}{v
  > 0}$ means $\msg{\reqOp{write}{$\iota$,v}}{v > 0}$. We omit
quantifiers with equality constraints, e.g.,
$\msg{\reqOp{write}{$\iota$,v}}{\iota = 1 \land v = 3}$, as
$\reqOp{write}{1\;3}$; we also omit the top ($\top$) qualifier
  in symbolic events and types.

\subsection{Typing Rules}

\begin{figure}[!t]
{\footnotesize
{\footnotesize
\begin{flalign*}
 &\text{\textbf{Auxiliary Typing}} & & \fbox{$\Gamma \vdash \phi \quad \Gamma \vdash A \subseteq A \quad \Gamma \vdash \tau <: \tau$} &\text{\textbf{Typing}} &  & \fbox{$\Gamma\vdash \eff{op}: t \quad \Gamma\vdash v : t \quad \Gamma \vdash e: \tau$}
\end{flalign*}
}
\\ \
\mprooftr{30mm}
{$\Gamma\wfvdash \tau[x\mapsto v]$}
{SubG}
{$\Gamma \vdash x{:}b\garr \tau <: \tau[x\mapsto v]$}
\;
\mprooftr{41mm}
{$\Gamma \vdash t_1 <: t_2$ \quad
$\Gamma,x{:}t_1 \vdash F_2 \subseteq F_1$ \quad
$\Gamma,x{:}t_1 \vdash H_1 \subseteq H_2$
 }
{SubHF}
{$\Gamma \vdash \uhat{H_1}{x{:}t_1}{F_1} <: \uhat{H_2}{x{:}t_2}{F_2}$}
\;
\mprooftr{30mm}
{$\Gamma \vdash \eff{op} : \overline{x{:}t_x}\sarr\uhat{H}{t}{F}$ \quad
$\Gamma,\overline{x{:}t_x} \vdash H' \subseteq H$}
{TOpHis}
{$\Gamma\vdash \eff{op} : \overline{x{:}t_x}\sarr\uhat{H'}{t}{F}$}
% \quad
% \mprooftr{30mm}
% {$\Gamma\vdash e: \tau$ \quad
% $\Gamma \vdash \tau <: \tau'$\quad $\Gamma \vdash \tau' <: \tau$}
% {TEq}
% {$\Gamma\vdash e: \tau'$}
\\[.8em]
\mprooftr{20mm}
{$\Gamma\vdash v : t$}
{TRet}
{$\Gamma\vdash v : \uhat{H}{t}{\epsilon}$}
\quad
\mprooftr{15mm}
{$\Delta(\eff{op}) = t$}
{TOpCtx}
{$\Gamma\vdash \eff{op} : t$}
\quad
\mprooftr{58mm}
{
$\Gamma \vdash \eff{op} : \overline{x_i{:}t_i}\sarr \uhat{H}{t}{\msgA{op}{\overline{x_i}} \seqA F}$ \quad
$\forall i. \;\Gamma \vdash v_i : t_i$}
{TEffOp}
{$\Gamma\vdash \eff{op}{\;\overline{v_i}} : \uhat{H}{t}{\msgA{op}{\overline{v_i}} \seqA F}$}
\\[.8em]
\mprooftr{56mm}
{
$\Gamma \vdash \eff{op} : \overline{x_i{:}t_i}\sarr \uhat{H}{t}{\msgA{op}{\overline{x_i}} \seqA F}$ \;
$\forall i. \;\Gamma \vdash v_i : t_i$}
{TReqOp}
{$\Gamma\vdash \reqOp{op}{\;\ensuremath{\overline{v_i}}} : \uhat{H}{\Code{id}}{\reqEvt{op}{\ensuremath{\overline{v_i}}} \seqA F}$}
% \mprooftr{62mm}
% {
% $\Gamma\vdash \eff{op} : \overline{x_i{:}t_i}\sarr \uhat{H}{\Unit}{\msgA{op}{\overline{x_i}} \seqA F}$ \quad
% $\forall i. \;\Gamma \vdash v_i : \widehat{t_i}$ \quad
% $\Gamma\vdash e : \uhat{H \seqA (\msgA{op}{\overline{v_i}} \land \msg{op}{\phi}) \seqA F}{t}{F'}$}
% {TGen}
% {$\Gamma\vdash \zgen{op}{\overline{v_i}}{e} : \uhat{H}{t}{\msg{op}{\phi} \seqA F \seqA F'}$}
\;
\mprooftr{51mm}
{
$\Gamma \vdash \eff{op} : \overline{x_i{:}t_i}\sarr \uhat{H}{t}{\msgA{op}{\overline{x_i}} \seqA F}$ \;
$\Gamma \vdash v : \Code{id}$}
{TRspOp}
{$\Gamma\vdash \respOpAlt{op}{\ensuremath{v}} : \uhat{H}{t}{\respEvt{op}{\ensuremath{\overline{v_i}}} \seqA F}$}
\\[.8em]
\mprooftr{35mm}
{
$\Gamma\vdash e_1 : \tau_1$ \quad $\Gamma\vdash e_2 : \tau_2$}
{TChoice}
{$\Gamma\vdash e_1 \oplus e_2 : \tau_1 \interty \tau_2$}
\quad
\mprooftr{66mm}{
  $\Gamma\vdash e_1 : \uhat{H}{x{:}t_x}{F_1}$ \quad
  $\Gamma, x{:}t_x\vdash e_2 : \uhat{H \seqA F_1}{t}{F_2}$}
{TLet}
{
  $\Gamma\vdash \zlet{x}{e_1}{e_2} : \uhat{H}{t}{F_1 \seqA F_2}$
}
\\[.8em]
\mprooftr{22mm}{
 $\Gamma, x{:}t_x\vdash e : \tau$
}{TFun}{
 $\Gamma\vdash \zzlam{x}{e} : x{:}t_x\sarr \tau$}
\
\mprooftr{47mm}{
}{TFixBase}{
 $\Gamma\vdash \zzfix{f}{x}{e} : \overline{x{:}s}\sarr\uhat{H}{\urt{b}{\bot}}{\emptyset}$}
\
\mprooftr{27mm}{
 $\Gamma\vdash \zzfix{f}{x}{e} : t$ \quad
 $\Gamma, f{:}t\vdash \zzlam{x}{e} : t'$
}{TFixInd}{
 $\Gamma\vdash \zzfix{f}{x}{e} : t'$}
% \quad
% \mprooftr{60mm}{
%  $\Gamma\vdash v_1 : y{:}t_y\sarr\uhat{H}{A}$ \quad
%  $\Gamma\vdash v_2 : t_y$ \quad
%  $\Gamma; \Delta' \vdash e' : \uhat{H \seqA A}{A'}$
% }{TApp}{
%  $\Gamma\vdash \zlet{x}{v_1\;v_2}{e} : \uhat{H}{A \seqA A'}$}
}
\vspace*{-.1in}
\caption{Selected typing rules. \BD{Does it make sense to not tag the
    future automata in \textsc{TRspOp} with a response tag? I think
    that should let history automata treat 'completed' asynchronous
    effects and synchronous effects uniformly; otherwise we might need
    to duplicate signatures.}}
\label{fig:selected-typing-rules}
\vspace*{-.15in}
\end{figure}

$\DSL{}$ is equipped with a pair of typing judgments on pure and
effectful expressions, $\Gamma \vdash v : t$ and $\Gamma\vdash e:
\uhat{H}{x{:}t}{F}$, respectively. Both judgments rely on a type
context, $\Gamma$, that maps from variables to pure refinement types
(i.e., $t$). As in prior work~\cite{ZMD+23}, typing contexts are not
allowed to contain \tyName{}s--- doing so breaks several structural
properties (e.g., weakening) that are used to prove type safety. Both
rules are parameterized over an additional context $\Delta$, which
provides types for pure and effectful operators.

\paragraph{Auxiliary typing relations}
Our type system also relies on three additional auxiliary
relations.\footnote{The full definition of these relations is included
  in the appendix.} The first is a well-formedness relation $\Gamma
\wellfoundedvdash \tau$ which ensures, among other things, that all
qualifiers appearing in a type $\tau$ are closed under the current
typing context $\Gamma$; the return and argument types of functions
are only refined with an under- and overapproximate qualifiers,
respectively; that the verification conditions generated by the
synthesis algorithm presented in the next section can be encoded as
effectively propositional (EPR) sentences~\cite{Ramsey1987}, which can
be efficiently handled by an off-the-shelf theorem prover such as
Z3~\cite{de2008z3}; and that the history trace of a $\resp$ operation
contains a matching $\req$ with the same id.

The type system also uses a mostly standard semantic subtyping
relation; \autoref{fig:selected-typing-rules} presents the key
subtyping rule for \tyName{}s.
The rule \textsc{SubG} allows to instantiate the ghost varaible in a \tyName{} into well-formed types. The rule \textsc{SubHF} uses on an
auxiliary inclusion relation on SREs $\Gamma \vdash A \subseteq A$ to check relate the history and future regexes of two \tyName{}s under the current type context $\Gamma$.
Notably, unlike (overapproximating)
trace-based refinement type systems~\cite{ZYDJ24}, this subtyping
relation is covariant in the history regex and contravariant in the
future regex, yielding behavior similar to Incorrectness Logic with
the reverse principle of consequence~\cite{OHP19}.  Since the future
automata in a \tyName{} can refer to the pure value it returns, the
inclusion check between the future automata extends the typing context
with a binding for the return value.

A subset of our typing rules is shown in
\autoref{fig:selected-typing-rules}.\footnote{The complete set of
  typing rules is included in the supplemental material.}  All our
rules assume any types they use are well-formed, so we elide the
corresponding well-formedness judgments from their premises.
The \textsc{TRet} rule requires the future regex of a pure term
($\epsilon$) to only accept the empty trace (i.e., $[]$). 
The operator type is retrieved from the operator context $\Delta$ using the \textsc{TOpCtx} rule.
For a given event, there may be multiple reachable history traces (underapproximate preconditions) that can be allowed by an operator, which cannot be determined by the operator type alone.
Therefore, we expect that the type provided by the operator context is the most general one, i.e., it includes all reachable history traces.
The rule \textsc{TOpHis} can then freely choose a subset of the reachable history traces to align with the actual calling context.
The application rule \textsc{TEffOp} requires the arguments of an effectful operator to be well-typed against the parameter types, and the \textsc{TRspOp} rule requires the response of an effectful operator to be well-typed against the return type.

\begin{example}[Operator Application Typing]
  We present a typing derivation for a simple generator that performs the $\eff{push}$ and $\eff{pop}$ operators in
  Example~\ref{ex:stack-example}. That is, we need to type check the following expression:
  \vspace{-.4cm} \par\nobreak {\small %
    \begin{align*}
      &\emptyset \vdash \zlet{x}{\mykw{assume}\;0 < x}{\eff{push}\;x; \eff{pop}} 
      \\&\qquad : \uhat{\msgAGray{pushI}{0\;0}\msgAGray{popI}{0\;0}}{x{:}\urt{\Int}{0 < \vnu}}{\msgA{push}{x}\msgAGray{pushI}{1\;x}\msgA{pop}{x}\msgAGray{popI}{1\;x}}
    \end{align*}}%
  \noindent 
  The generator assume the stask is empty (i.e., the history regex contains $\effGray{pushI}$ and $\effGray{popI}$ events with $0$ index and dummy value $0$) and push a postive number $y$ and expect to pop the same value.
  The $\mykw{assume}\;\true$ can guarantee generates all positive integers, thus we can type check two effectful operation applications via the following judgements:
  \vspace{-.4cm} \par\nobreak {\small %
    \begin{flalign*}
      &x{:}\urt{\Int}{0 < \vnu} \vdash \eff{push}\;x; \eff{pop} :  \uhat{\msgAGray{pushI}{0\;0}\msgAGray{popI}{0\;0}}{x{:}\urt{\Int}{0 < \vnu}}{\msgA{push}{x}\msgAGray{pushI}{1\;x}\msgA{pop}{x}\msgAGray{popI}{1\;x}}
    \end{flalign*}}%
  \noindent
  According to the \textsc{TEffOp} rule, we first retrieve the type of the $\eff{push}$ operators from the operator context $\Delta$:
  \vspace{-.4cm} \par\nobreak {\small %
    \begin{flalign*}
      &... \vdash \eff{push} : x{:}\urt{\Int}{y < \vnu} \sarr
      \uhat{\Code{Last}(\msgAGray{pushI}{n\;y})}{\Unit}{\msgA{push}{x}\msgAGray{pushI}{n{+}1\;x}}  &&\text{(by \textsc{TOpCtx})}
      \\&... \vdash \eff{push} : x{:}\urt{\Int}{0 < \vnu} \sarr
      \uhat{\Code{Last}(\msgAGray{pushI}{0\;0})}{\Unit}{\msgA{push}{x}\msgAGray{pushI}{1\;x}} &&\text{(by \textsc{SubG})}
      \\&... \vdash \eff{push} : x{:}\urt{\Int}{0 < \vnu} \sarr
      \uhat{\msgAGray{pushI}{0\;0}\msgAGray{popI}{0\;0}}{\Unit}{\msgA{push}{x}\msgAGray{pushI}{1\;x}} &&\text{(by \textsc{TOpHis})}
    \end{flalign*}}%
  \noindent 
  The subtyping rule allow us to instantiate the ghost varaible (e.g., $n \mapsto 0$ and $y \mapsto 0$), moreover, the rule \textsc{TOpHis} choose one reachable history trace ($\msgAGray{pushI}{0\;0}\msgAGray{popI}{0\;0}$) to align with the actual calling context. Then, the future regex $\msgA{push}{x}\msgAGray{pushI}{1\;x}$ is exactly aligned with the prefix of the prefix of the future regex of the target type.
  The rest of program should be typed via the following judgements:
  \vspace{-.4cm} \par\nobreak {\small %
    \begin{align*}
      &x{:}\urt{\Int}{0 < \vnu}\vdash \eff{pop} : \uhat{\msgAGray{pushI}{0\;0}\msgAGray{popI}{0\;0}\msgA{push}{x}\msgAGray{pushI}{1\;x}}{x{:}\urt{\Int}{0 < \vnu}}{\msgA{pop}{x}\msgAGray{popI}{1\;x}} &&
    \end{align*}}%
  \noindent
  Similarly, we can instantiate the ghost varaible $m$ as $0$ and $n$ as $1$ via rule \textsc{SubG}, align the return type via rule \textsc{SubHF} and refine the history regex via the rule \textsc{TOpHis}.
  \vspace{-.4cm} \par\nobreak {\footnotesize %
    \begin{flalign*}
      &... \vdash \eff{pop} : \uhat{\Code{Last}(\msgAGray{popI}{m\;\_}) \land
      \Code{Last}(\msgAGray{pushI}{n\;\_}) \land
      (\allA\msgAGray{pushI}{(n{-}m)\;x}\allA)}{x{:}\Int}{\msgA{pop}{x}\msgAGray{popI}{m{+}1\;x}} &&\text{(\textsc{TOpCtx})}
      \\&... \vdash \eff{pop} : \uhat{\Code{Last}(\msgAGray{popI}{0\;\_}) \land
      \Code{Last}(\msgAGray{pushI}{1\;\_}) \land
      (\allA\msgAGray{pushI}{1\;x}\allA)}{x{:}\Int}{\msgA{pop}{x}\msgAGray{popI}{1\;x}} &&\text{(\textsc{SubG})}
      \\&... \vdash \eff{pop} : \uhat{\Code{Last}(\msgAGray{popI}{0\;\_}) \land
      \Code{Last}(\msgAGray{pushI}{1\;\_}) \land
      (\allA\msgAGray{pushI}{1\;x}\allA)}{x{:}\urt{\Int}{0 < \vnu}}{\msgA{pop}{x}\msgAGray{popI}{1\;x}} &&\text{(\textsc{SubHF})}
      \\&... \vdash \eff{pop} : \uhat{\msgAGray{pushI}{0\;0}\msgAGray{popI}{0\;0}\msgA{push}{x}\msgAGray{pushI}{1\;x}}{x{:}\urt{\Int}{0 < \vnu}}{\msgA{pop}{x}\msgAGray{popI}{1\;x}} &&\text{(\textsc{TOpHis})}
    \end{flalign*}}%
  \noindent
\end{example}

% \begin{example}[Operator Type Context] The type of effect operator
%   $\eff{readRsp}$ in Example~\ref{ex:read-atomicity} can provided by
%   the following operator context $\Delta$: %
%   \vspace{-.4cm} \par\nobreak {\small %
%     \begin{align*}
%       \Delta \equiv \{&(\eff{readRsp}, \iota{:}\Int\sarr \uhat{\msgA{writeRsp}{\iota\;x}\msgA{readReq}{\iota}}{x{:}\urt{\Int}{\top}}{\msgA{readRsp}{\iota\;x}}),\\
%                       &(\eff{readRsp}, \iota{:}\Int\sarr
%                         \uhat{\msgA{writeRsp}{\iota\;x}\msgA{readReq}{\iota}\msgC{writeRsp}}{x{:}\urt{\Int}{\top}}{\msgA{readRsp}{\iota\;x}}),
%                         ... \}
%     \end{align*}}
% \end{example}

\BD{I think we need a couple of examples of how we use the subtyping
  rules to specialize / unify the history automata of a effectful
  operator with the current context, and what differentiates how this
  type system does so versus \textsc{Hat}s, analgously to what we did
  for the subtyping relation in the coverage type work. }

\BD{What does the typing rule for assert statements look like? Is it
  more interesting in the underapproximate setting, in that a
  well-typed assert only needs to be guaranteed to succeed in at least
  one execution?}

\SJ{Why isn't the type given in Sec. 2 sufficient?}

% The typing rules effectful operations, \textsc{TGen} and
% \textsc{TObs}, both extract the operator type
% $\uhat{H}{t}{\msg{op}{\overline{x}}\seqA F}$ from operator context
% $\Delta$ \BD{The current rules do not mention $\Delta$, need to fix};
% and require that it aligns with the \tyName{} of the expression that
% the operation is being performed in:%
% The rules for performing events, \textsc{TGen} and \textsc{TObs}, both extract
% the type of the corresponding operator from $\Delta$,
% $\rg{H}{\msg{op}{\phi}}{A \seqA F}$, and require that it aligns
% with the \tyName{} of the expression that the operation is being
% performed in:%
% \vspace{-.4cm} \par\nobreak {\footnotesize %
%   \begin{align*}
%     \text{\textcolor{gray}{(term's \tyName)}} \qquad
%     \effLR{\underbrace{ H}_\texttt{history}}\;t\;
%     \effLR{
%     \;\;\underbrace{\msg{op}{\overline{x}}}_\texttt{trace produced by $\eff{op}$}
%     \;\;\underbrace{F }_\texttt{ghost events appended to $\eff{op}$}
%     \;\;\underbrace{F'}_\texttt{trace produced by $e$}
%     }
% \end{align*}}%
% \SJ{I don't follow what it means to append a ghost event to a normal
%   event.}
% \noindent For simplicity, we only allow the ghost events to be
% appended to a normal event.  Same as prior work~\cite{CoverageType},
% the arguments should be well-typed against the parameter types fliped
% from over- to under-approximation (i.e., $\widehat{t_i}$).  To type
% the rest of the expression, the assumption of \tyName of operator
% (i.e., $\msg{op}{\overline{x}}$) should be aligned with the
% expectation of term's type (i.e., $\msg{op}{\phi}$), then both rules
% move the symbolic event
% $(\msg{op}{\overline{x}} \land \msg{op}{\phi})\seqA F$ from the future
% regex to the tail of the history regex. The subsumption rule
% \textsc{TSub} and the equivalence rule \textsc{TEq} allows us to
% change the shape of a type that a term is being typed against.  

The nondeterministic choice operator is typed using the \textsc{TChoice}
rule which merges \tyName{} of the two branches.  Although the typing
rule for function (\textsc{TFun}) is standard, unlike standard
refinement type systems, underapproximate reasoning does not maintain
an inductive invariant for recursive functions. Inspired by
Underapproximate Logic, we type recursive functions in an inductive
manner for instead. The \textsc{FixBase} rule types any recursive
function with a \tyName{} lacking reachability constraints (i.e.,
$\uhat{H}{\urt{b}{\bot}}{\emptyset}$). In contrast, the
\textsc{FixInd} rule permits typing a recursive function with type
$t'$ by requiring it to have another type $t$, and by checking that
its body can be assigned type $t'$ under the assumption that the
function itself has type $t$.  \SJ{This description is confusing -
  what is the intuition used in the IL treatment of recursion that
  we can refer to and apply here?}

% The nondeterministic loop aligns with the Kleene star in SRE, and thus can be straightforwardly typed by the \textsc{TLoop} rule.
% Since the loop body $e$ may be executed an unbounded number of times, the typing judgment for $e$ requires that an unbounded number of iterations may have already been performed ($H \seqA \textcolor{orange}{A^*}$) and may occur in the future ($\textcolor{orange}{A^*} \seqA A' \seqA F$). Additionally, the input and output capability contexts must be identical, ensuring that $e$ replenishes any capabilities it consumes.

% of observing effect operator $\eff{op}$ from the typing
% context with the help of effect operator typing
% $\textsc{TEOp}$. Besides requiring the arguments ($\overline{v_i}$) is
% well-typed against parameter types ($\overline{t_i}$) as application
% rules in standard refinement typing, it additionally

\subsection{Metatheory}
\label{sec:metatheory}
% \begin{figure}[h!]
%   \vspace*{-.15in}

%   \vspace*{-.15in}
%   \caption{Selected type denotation.}
%   \label{fig:selected-denotation}
%   \vspace*{-.15in}
%   \end{figure}

\begin{figure}[t!]
\vspace*{-0.1in}
{\footnotesize
\begin{flalign*}
  &\text{\textbf{Operator Semantics }}\
  \fbox{$\alpha \vDash m \Downarrow m$} &
  \text{\textbf{Operational Semantics }}\
  \fbox{$
  % \phi \Downarrow c \quad
  \steptr{\alpha}{(\buf, e)}{\alpha}{(\buf, e)}$}
\end{flalign*}
\begin{prooftree}
  \hypo{\beta' = \buf \cup \{\zevent{op}{\overline{c}}\}}
  \infer1[\textsc{\small StGen}]{
    \steptr{\alpha}{(\buf, \zgen{op}{\overline{c}}{e})}{[\zevent{op}{\overline{c}}]}{(\beta', e)}
  }
\end{prooftree}
\quad
\begin{prooftree}
  \mhypo{50mm}
  {$\buf = \{ m \} \cup \buf' $ \quad
    $\alpha \vDash m \Downarrow \zevent{op}{\overline{c},c'} $
  }
    \infer1[\textsc{\small StObs}]{
      \steptr{\alpha}{(\buf,
        \zobs{x}{op}{\overline{c}}{e})}{[\zevent{op}{\overline{c},c'}]}{(\buf', e[x \mapsto c'])}
    }
  \end{prooftree}
% \\[0.8em]
% \mprooftr{58mm}{
% $\phi\msubst{x}{c}\Downarrow \top$
% }{StAssume}{
% $\steptr{\alpha}{(\buf, \zassume{\overline{x}}{\phi}{e})}{[]}{(\buf, e\msubst{x}{c})}$
% }
% \quad
% \mprooftr{44mm}{$\phi \Downarrow \top$}{StAssert}{
% $\steptr{\alpha}{(\buf, \zassert{\phi}{e})}{[]}{(\buf, e)}$
% }
}
\vspace*{-.05in}
  \caption{Selected Operational Semantics.}
  \label{fig:selected-semantics}
\vspace*{-.2in}
\end{figure}

\paragraph{Operational Semantics} Events in $\DSL$ are operations
applied to concrete values $(\zevent{op}{\overline{c}})$. Evaluating a
$\DSL$ program depends on an context \emph{trace}, i.e., a sequence of
events, and an \emph{observation capability}, i.e., a multiset of
generable events that are waiting to be handled by observable
operators. Each evaluation step produces an output trace and an
updated capability. Traces are equipped with the standard list
operations (i.e., cons ${::}$ and concatenation $\listconcat{}$).  The
operational semantics of $\DSL$ are defined by the small-step
reduction relation:
$\steptr{\alpha}{(\beta, e)}{~\alpha'}{(\beta', e')}$.  This judgment
is read as: ``under the history trace $\alpha$ and observation
capability $\beta$, $e$ steps to $e'$, emitting the trace $\alpha'$
and remaining the capability $\beta'$.'' Intuitively, the history
trace $\alpha$ represents the sequence of events visible to a
operator, thereby determining its response; the capability $\beta$
contains events that are able to be observed in the future. The
semantics uses an auxiliary judgement,
$\alpha \vDash m \Downarrow m'$, that specifies the observable event
$m'$ that will be generated after handling a generative event $m$.

\autoref{fig:selected-semantics} presents the key rules of $\DSL$'s
semantics,\footnote{The remaining rules are completely standard and
  provided in the supplemental material.} which reflect the
``consume-and-produce'' semantics of observation capability.  The rule
for observable events (\textsc{StObs}) removes an event $m$ from the
capability that can produce the observable operation $\eff{op}$,
evaluates it under the current context, and substitutes the observed
value $c'$ (last payload of $m$) for binding variable $x$ in $e$, the
body of the $\mykw{let}$ expression.  In contrast, the rule for
generable events (\textsc{StGen}) makes no assumptions on the
observation capability, since these events can be directly performed
by the generator. It simply add itself to the resulting observation
capability for future handling.

\begin{example}[Auxiliary Operator Semantics]\label{ex:aux-semantics}
  The auxiliary semantics of the $\eff{read}$ and $\eff{write}$ operators in Example~\ref{ex:read-atomicity} can be defined as follows:

  {\footnotesize
    \mprooftr{44mm}{}{OpWrite}{
    $\alpha \vDash \zevent{writeReq}{\iota,v} \Downarrow \zevent{writeRsp}{\iota,v}$
    }
    \quad
    \mprooftr{62mm}{$\alpha_1 = \alpha_{11} \listconcat \zevent{writeRsp}{\iota',v} \listconcat \alpha_{12}$ \quad $\eff{readRsp}\not\in \alpha_{12}$}{OpRead}{
    $\alpha_1\listconcat\zevent{readReq}{\iota}\listconcat\alpha_2\vDash \zevent{readReq}{\iota} \Downarrow \zevent{readRsp}{\iota,v}$
    }
    \\[-0.1em]
    }\noindent

  \noindent
  The database responds to a $\eff{writeReq}$ event by emitting a
  $\eff{writeRsp}$ event containing the identical payload. To ensure
  read atomicity, handling a $\eff{readReq}$ event requires
  identifying the corresponding history prefix $\alpha_1$ at the point
  of the read request happens, and producing a $\eff{readRsp}$ event
  whose value $v$ matches the most recently written value for that key
  in $\alpha_1$.
\end{example}

\begin{example}[Read Atomicity]\label{ex:semantics-async}
  The following derivation illustrates the execution of a single iteration of the program from Example~\ref{ex:read-atomicity}, evaluated under the auxiliary semantics for $\eff{writeReq}$ and $\eff{readReq}$ specified in Example~\ref{ex:aux-semantics}, which enforce read atomicity. For this iteration, we assign transaction identifiers $\Code{tid1}=1$ and $\Code{tid2}=2$, and two random generated integers are $7$ and $8$:
    {\footnotesize
      \begin{flalign*}
        &[] \vDash
        (\emptyset, \eff{writeReq}\;1\;7;\zobs{\_}{writeRsp}{1}{\eff{readReq}\;2; \textit{line 5 - 6 of code in Example~\ref{ex:read-atomicity}}}) &&
        \\[-0.2em]
        &\;\;\xhookrightarrow{[\zevent{writeReq}{1,7}]}{(\{ \zevent{writeReq}{1,7} \}, \zobs{\_}{writeRsp}{1}{\eff{readReq}\;2; \textit{line 5 - 6 of code in Example~\ref{ex:read-atomicity}}})}  \tag{\textsc{StGen}} &&
        \\[-0.4em]
        &\;\;\xhookrightarrow{[\textcolor{blue}{\zevent{writeRsp}{1,8}}]}{(\emptyset, \eff{readReq}\;2; \eff{writeReq}\;1\;7;\zobs{\_}{writeRsp}{1}{\textit{line 6 of code in Example~\ref{ex:read-atomicity}}})}  \tag{\textsc{StObs}} &&
        \\[-0.4em]
        &\;\;\xhookrightarrow{[\textcolor{orange}{\zevent{readReq}{2}]}}{(\{\zevent{readReq}{2}\}, \eff{writeReq}\;1\;8;\zobs{\_}{writeRsp}{1}{\textit{line 6 of code in Example~\ref{ex:read-atomicity}}})}  \tag{\textsc{StGen}} &&
        \\[-0.4em]
        &\;\;\xhookrightarrow{[\zevent{writeReq}{1,8}]}{(\{\zevent{readReq}{2}, \zevent{writeRsp}{1,8}\}, \zobs{\_}{writeRsp}{1}{\textit{line 6 of code in Example~\ref{ex:read-atomicity}}})}  \tag{\textsc{StGen}} &&
        \\[-0.4em]
        &\;\;\xhookrightarrow{[\zevent{writeRsp}{1,8}]}{(\{\zevent{readReq}{2}\}, \zobs{y}{readRsp}{2}{\sassert{y == 7}; () \oplus \Code{transaction\_gen\;1\;2}})}  \tag{\textsc{StObs}} &&
        \\[-0.4em]
        &\;\;\xhookrightarrow{[\zevent{readRsp}{2,7}]}{(\emptyset, \sassert{7 == 7}; () \oplus \Code{transaction\_gen\;1\;2})}  \tag{\textsc{StObs}} &&
      \end{flalign*}
    }\noindent
  In the final step, the program executes $\mykw{obs}\;\eff{readRsp}\;2$, producing a read value $y=7$, which matches the value of the most recent $\eff{writeRsp}$ (blue) at the point where the $\eff{readReq}$ event (orange) occurred. Thus, read atomicity is preserved and execution can continue to generate further transactions.
\end{example}

\paragraph{Type denotations} Similar to other refinement type systems~\cite{JV21},
types in $\DSL$ are denoted as their inhabitants (i.e., $\denot{t}$ and $\denot{\tau}$).  The capability context is denoted as  observation capabilitys (i.e., $\denot{\Theta}$), while the type context $\Gamma$ is denoted as \emph{substitution} $\sigma$ (i.e., $[x_1\mapsto v_1, x_2\mapsto v_2]$) that provides the assignments for binding variables in $\Gamma$. Moreover, the denotation (denoted language) of SREs is the set of traces. Then, regex
inclusion under a type context is defined as $\Gamma \vdash A \subseteq
A' \doteq \forall \sigma \in \denot{\Gamma}. \denot{\sigma(A)}
\subseteq \denot{\sigma(B)}$.
The type denotation of $\uhat{H}{x{:}\urt{b}{\phi}}{F}$ is defined as follows:
% {\small\begin{align*}
%   &\denot{\rg{H}{A}{F}}_{(\Theta, \Theta')} \doteq \{(e,  e_f) ~|~ \emptyset \basicvdash e : \Code{unit} \land
%   \\&\qquad \forall \alpha \in \denot{A}. \exists \alpha_h \in\denot{H}.\exists \alpha_f \in \denot{F}. \exists \beta \in \denot{\Theta}. \exists \beta' \in \denot{\Theta'}. \msteptr{\alpha_h}{(\beta, e;e_f)}{\alpha}{(\beta',  e_f)}\mharr{\alpha_f}{(\emptyset, ())} \}
% \end{align*}}\noindent
{\small\begin{align*}
  % &\denot{\uhat{H}{x{:}\urt{b}{\phi}}{F}} \doteq 
  % \\[-0.4em]& \quad 
  \{e ~|~ \emptyset \basicvdash e : s \land \forall v{:}b. \phi[v\mapsto v] \impl \forall \alpha_f \in \denot{F[x\mapsto v]}. \exists \alpha_h \in\denot{H[x\mapsto v]}. \exists \beta\;\beta'. \msteptr{\alpha_h}{(\beta, e)}{\alpha_f}{(\beta',  v)} \}
\end{align*}}\noindent
A term $e$ inhabits $\uhat{H}{x{:}\urt{b}{\phi}}{F}$ iff, for trace $\alpha_f$ accepted by $F$, there exists history trace $\alpha_h$ accept by $H$, such that the execution of $e$ under history trace produces the trace $\alpha_f$.
\footnote{The details of these definition can be found in our \techreport{}.}
Our formulation constitutes a trace-based analogue of incorrectness logic~\cite{OHP19}, in which effects are modeled as program state transitions:
\begin{align*}
  [P]\; e\; [Q] \doteq \forall s_\Code{post} \in \denot{Q}. \exists s_\Code{pre} \in \denot{P}. (s_\Code{pre}, e) \Downarrow s_\Code{post}
\end{align*}
The distinction is that we characterize $F$ as the set of newly produced traces rather than the full post-execution trace, reflecting our trace-based semantics. This avoids redundant prefixes inherited from the execution history.

% In contrast to incorrectness logic~\cite{OHP19}, which explicitly models effects via program states, our latent effect semantics require executions to terminate with an empty capability context. Permitting termination with a non-empty capability (e.g., $\{\eff{readRsp}\}$) would leave latent effects from asynchronous actions (such as $\eff{readReq}$) unobserved. Thus, we explicitly denote the \tyName with term and its continuation $e_f$ to capture the entire execution. The $\beta$ and $\beta'$ are the capability before and after execution of $e$; since it is always valid to add unused capability to both sides, thus this denotation is parameterized by capability context $\Theta_A$ and $\Theta_F$.

% Underapproximating specifications, as in our approach, must exclude unreachable states or unrealizable traces, since $A$ cannot include them. To improve expressiveness and usability, our tool allows users to specify an overapproximated style global property $P$ and verifies that $A \subseteq P$.

%  A term $e$ inhabits $\rg{H}{A}{F}$ iff, for all realizable (i.e., producible by the execution of some program) history traces $\alpha_h$ accept by $H$, the current trace $\alpha$ produced by $e$ is accepted by $A$, and all future traces $\alpha_f$ accepted by $F$ are also realizable.\footnote{The details of these definition can be found in our \techreport{}.}

% \paragraph{Operator context}
% We expect the trace produced by our test generator to be consistent with the operator context $\Delta$; i.e., each event in the trace must be consistent with corresponding $\tyName$ in $\Delta$.

% \begin{definition}[Trace consistent with operator context]
%   \label{lemma:trace-conistency}
%   A trace $\alpha$ is consistent with a operator context $\Delta$
%   (written $\traceDelta{\Delta}{\alpha}$) iff for every
%   $\zevent{op}{\overline{c_i}}$ in $\alpha$ (i.e., $\alpha = \alpha_h
%   \listconcat [\zevent{op}{\overline{c_i}}] \listconcat \alpha_f$),
%   there exists $\overline{x{:}t} \sarr \rg{H}{\msg{op}{\phi}}{F}$ and
%   capability $\Theta$ in $\Delta$ such that the following hold: %
%   \vspace{-.4cm} \par\nobreak {\small %
%     \begin{align*}
%     \alpha_h \in \denot{H\msubst{x_i}{c_i}} \land \alpha_f \in \denot{F\msubst{x_i}{c_i}} \land \Theta \subseteq \{ \eff{op} ~|~ \zevent{op}{\overline{c}} \in \alpha_f\} \land \phi\msubst{x_i}{c_i} \land \forall i. c_i \in \denotation{t_i}
% \end{align*}}
% \end{definition}

\begin{theorem}[Fundamental Theorem]\label{theorem:fundamental}A well-typed term, i.e., \(\Gamma\vdash e : \uhat{H}{t}{F}\), generates traces consistent with the \tyName: $\forall \sigma \in \denotation{\Gamma}. \sigma(e) \in \denotation{\sigma(\uhat{H}{t}{F})}$.\footnote{The proofs of all theorems in this paper are provided in the \techreport{}.}
\end{theorem}

\begin{corollary}[Type Soundness]\label{theorem:sound} A generator \(e\) that satisfies \(\emptyset \vdash e : \uhat{\epsilon}{\Unit}{F}\) must produce all traces accepted by $F$, i.e., $\alpha \in \denot{F}. \exists \beta.\msteptr{[]}{(\emptyset, e)}{\alpha}{(\beta, ())}$.
\end{corollary}
